\documentclass[level=3,teacher]{kidscode}
%% Terminal and text-file listings: no Python keyword coloring.
\lstdefinestyle{kcshell}{basicstyle=\ttfamily\small\setstretch{1.0},language={},
  columns=fullflexible,keepspaces=true,breaklines=true,numbers=none,upquote=true,
  showstringspaces=false}
\begin{document}
\kctitle{Teacher Guide}{Everything the parent-teacher needs for Level 3: Advanced}

\begin{tip}[How to use this guide]
Read sections 1--6 once before the first session (about 45 minutes, spread over a couple of
evenings). After that you only need the weekly lesson plan (10 minutes of reading each week)
and this guide's \textbf{Troubleshooting} section when something goes wrong.
\end{tip}

\setcounter{tocdepth}{1}
\tableofcontents
\clearpage

%% =====================================================================
\section{1. Welcome to Level 3}

Your coders have come a long way: from Scratch blocks in Level 1 to Python games with
functions, lists, dictionaries, files and their own pytest tests in Level 2. Level 3 is where
they start working like real software makers. They will:

\begin{itemize}[itemsep=2pt]
  \item save their work with \textbf{git}, the same save-point tool professional programmers use;
  \item design programs out of \textbf{classes and objects} (characters, mazes, galleries);
  \item meet famous \textbf{algorithms} (binary search, sorting, breadth-first search, recursion)
        and talk about how fast they are;
  \item clean and chart \textbf{real data}, and build two tiny \textbf{AI} programs to see how
        machines learn from examples, and where they go wrong;
  \item plan, build, test, document and present a \textbf{larger project} of their own.
\end{itemize}

\subsection{What changes at Level 3}

\begin{tabularx}{\linewidth}{@{}>{\bfseries\raggedright\arraybackslash}p{3.2cm}X@{}}
\toprule
More independence & You explain less and ask more. Students read the workbook chapter, try
first, and ask a partner before asking you (\enquote{ask three, then me}). Your job shifts from
explaining to coaching: \enquote{What did you expect? What happened? What will you try next?}\\ \midrule
Projects, not exercises & From week 9 the group works on one larger project each (or per pair).
Weekly exercises become building blocks they can reuse.\\ \midrule
Terminal & Students type commands in the Mac Terminal: \texttt{git} for save points and
\texttt{python3 -m pytest} for tests. You will learn the five commands they use; the cheat
sheet lists them.\\ \midrule
Tests and history matter & A project is not only \enquote{does it run?} but also \enquote{does it
have tests?}, \enquote{can I see how it grew in the git log?} and \enquote{can you explain why you
designed it like that?}\\ \midrule
Bigger ideas & Efficiency, recursion, data and AI are ideas they will meet again for years. Aim
for a good intuition, not perfect definitions.\\
\bottomrule
\end{tabularx}

\begin{teachernote}[You do not need to be an expert]
Every lesson plan tells you exactly what to say, type and look for, and every program in the
answer key is tested. When a student asks something you do not know, the best answer is:
\enquote{Great question --- let's find out together.} Then try it in Thonny. Modeling curiosity
matters more than knowing everything.
\end{teachernote}

%% =====================================================================
\section{2. CS background for the parent}

Ten short ideas, in plain words. Each one is taught properly in its week; this is only so that
nothing in the lesson plans surprises you.

\begin{description}[leftmargin=0pt,itemsep=6pt,font=\bfseries]
\item[Class and object (weeks 2--3).] A \textbf{class} is a blueprint, like a cookie cutter. An
\textbf{object} is one thing made from it, like one cookie. \py|class Pet:| describes what every
pet has (\textbf{attributes}, such as \py|self.name| and \py|self.energy|) and can do
(\textbf{methods}, such as \py|play()|). \py|mo = Pet("Mo")| makes one pet. \py|self| simply
means \enquote{this particular object}.

\item[Inheritance and composition (week 3).] \textbf{Inheritance} is an \enquote{is a}
relationship: a \py|Knight| \emph{is a} \py|Character|, so \py|class Knight(Character):| gets all
of Character's code for free and only changes what is different (\textbf{overriding}).
\textbf{Composition} is a \enquote{has a} relationship: a \py|Game| \emph{has a} \py|Maze| and a
list of ghosts. Good designs use both.

\item[Big-O intuition: \enquote{how fast is fast?} (weeks 4--5).] We do not time programs with a
stopwatch; we count \textbf{steps} and ask how the count grows when the data gets bigger.
Checking every page of a 1000-page book one by one can take 1000 steps (\textbf{linear}). Opening
the middle, then the middle of the correct half, and so on, takes at most about 10 steps
(\textbf{binary search}, which needs sorted data). Doubling the book adds one step for binary
search but doubles the work for linear search. That is the whole idea; no formulas are needed.

\item[Sorting (week 5).] Putting items in order. Students act out \textbf{selection sort} (find the
smallest, swap it to the front, repeat) and \textbf{bubble sort} (swap neighbors that are in the
wrong order, repeat), then learn Python's built-in \py|sorted(items, key=...)|.

\item[Stack and queue (week 5).] A \textbf{stack} is a pile of plates: the last one put on is the
first taken off (an undo button). A \textbf{queue} is a lunch line: first in, first out (players
waiting for a turn). In Python a list works as a stack (\py|append|, \py|pop|) and
\py|collections.deque| works as a queue (\py|append|, \py|popleft|).

\item[Grid and recursion (week 7).] A \textbf{grid} is a list of lists: \py|grid[row][col]|. A
\textbf{recursive} function solves a problem by calling itself on a smaller version of the same
problem, like Russian dolls. It must have a \textbf{base case} where it stops; without one Python
gives up with a \texttt{RecursionError}. Fractal trees and flood fill (the paint bucket tool) are
recursive.

\item[Breadth-first search, BFS (week 7 extension and the final project).] How a computer finds the
shortest way through a maze. Imagine dropping ink at the start: it spreads one square in every
direction, then one more, ring after ring. The first time the ink reaches the goal, that is a
shortest path. A \textbf{queue} keeps the squares in the order the ink reached them. The Maze
Muncher reference project uses BFS for its chasing ghost.

\item[Data cleaning (week 8).] Real data is messy: empty rows, extra spaces, numbers stored as text.
\textbf{Cleaning} means fixing those before counting or charting. A \textbf{CSV} file is a
spreadsheet saved as plain text with commas between the columns.

\item[Nearest-neighbor classifier (week 9).] The simplest machine learning. You have labeled
examples (\textbf{training data}), each described by numbers (\textbf{features}, such as size and
number of spikes). To label a mystery item, find the example most similar to it (smallest
distance) and copy its label. If the examples are unfair or incomplete, the predictions are too:
that is \textbf{bias}. A \textbf{Markov text generator} (also week 9) writes new sentences by
remembering which word followed which in its training text --- it imitates without understanding,
a good first picture of how chatbots can sound confident and still be wrong.

\item[git (week 1 onwards).] A \textbf{save-point system for a folder}. \texttt{git add} chooses
which changes go into the next save point, \texttt{git commit -m "message"} makes the save point
(a \textbf{commit}), \texttt{git log} lists them, \texttt{git diff} shows what changed and
\texttt{git restore} throws away changes that were not committed. Everything stays on your Mac: no
account, no internet, no GitHub.
\end{description}

%% =====================================================================
\section{3. Level 3 goals}

Every lesson objective is tagged with one of these codes. CSTA codes are from the CSTA K--12
Computer Science Standards (revised 2017); AI4K12 refers to the \enquote{Five Big Ideas in AI}
(1 perception, 2 representation \& reasoning, 3 learning, 4 natural interaction, 5 societal impact).

\renewcommand{\arraystretch}{1.25}
\begin{tabularx}{\linewidth}{@{}>{\bfseries}p{1.3cm}X>{\raggedright\arraybackslash}p{3.9cm}@{}}
\toprule
\rowcolor{kcaccent!15} Code & By the end of Level 3 the student can\ldots & CSTA / other\\ \midrule
L3.G1 & design programs with classes, objects, inheritance and composition & 3A-AP-17, 3A-AP-18\\
L3.G2 & explain and compare searching/sorting algorithms and use recursion & 3A-AP-15, 1B-AP-08 (E), 3A-IC-26\\
L3.G3 & choose data structures: lists of objects, dicts, stacks, queues, 2D grids & 3A-AP-14, 3A-DA-10\\
L3.G4 & track their work with git (init, add, commit, log, diff, restore) & 3A-AP-22, 2-AP-19\\
L3.G5 & plan and build a larger program iteratively, using feedback & 3A-AP-16, 3A-AP-19, 3A-AP-23, 2-AP-18\\
L3.G6 & collect, clean, analyze and visualize data & 2-DA-08, 3A-DA-11, 3A-DA-12\\
L3.G7 & explain how simple AI learns from data and where it fails & AI4K12 Big Ideas 1--5, 3A-IC-24\\
L3.G8 & test and debug larger programs & 2-AP-17, 3A-AP-21\\
L3.G9 & reason about privacy, bias, ethics and licenses & 3A-IC-24, 3A-IC-25, 3A-IC-29, 3A-AP-20\\
\bottomrule
\end{tabularx}

%% =====================================================================
\section{4. The 12-week roadmap}

\begin{longtable}{@{}>{\bfseries}p{0.8cm}>{\raggedright\arraybackslash}p{3.6cm}>{\raggedright\arraybackslash}p{5.6cm}>{\raggedright\arraybackslash}p{\dimexpr\linewidth-10cm-6\tabcolsep}@{}}
\toprule
\rowcolor{kcaccent!15} Wk & Title & Big idea & What they build / do\\ \midrule\endhead
\bottomrule\endfoot
1 & Save Points: Git Basics & Git is a save-game system for code. & Placement quiz (15 min); first repository, first commits\\ \midrule
2 & Objects: Blueprints and Things & A class is a blueprint; objects are things built from it. & A virtual pet class\\ \midrule
3 & Families of Objects & Share code with inheritance; build with composition. & A family of game characters / shapes\\ \midrule
4 & Searching and How Fast Is Fast? & Smart algorithms do less work. & Linear vs binary search, counting steps\\ \midrule
5 & Sorting, Stacks and Queues & Order data; stacks undo, queues take turns. & Sorting a leaderboard, an undo stack, a turn queue\\ \midrule
6 & Review and Checkpoint & Connect the pieces. & Mid-level checkpoint (15 min); Thonny debugger; bug-hunt workshop\\ \midrule
7 & Grids, Recursion and Fractal Art & Functions that call themselves; worlds made of grids. & Fractal trees, flood fill on a grid\\ \midrule
8 & Data Detectives & Data tells stories when you clean and chart it. & Survey CSV $\rightarrow$ counts, averages, turtle bar chart\\ \midrule
9 & How Do Machines Learn? & AI finds patterns in examples --- and copies their mistakes. & Nearest-neighbor classifier, Markov text; project kickoff (pitch)\\ \midrule
10 & Project Studio 1 & Design first: classes, data, milestones. & CRC cards, class diagram, core classes + first tests\\ \midrule
11 & Project Studio 2 & Test, playtest, polish, document. & 5+ tests, playtest, README, polish\\ \midrule
12 & Showcase and Graduation & Share, explain, reflect, plan next steps. & Showcase, exit assessment (20 min), certificates\\ \midrule
B1 & Algorithm Olympics & Algorithms are recipes you can race. & Unplugged sorting networks, maze solving by hand\\ \midrule
B2 & Mini Game Jam & Small ideas, quick builds, many commits. & 90-second game ideas, free-build, git practice\\
\end{longtable}

%% =====================================================================
\section{5. Setup checklist (once, about 60 minutes)}

Do this on each Mac a week before session 1. Nothing here needs an Apple ID, a GitHub account or
any personal data. Internet is needed once, for the Apple download.

\subsection{5.1 git: Apple Command Line Tools}
\begin{checklist}
  \item Open \textbf{Terminal}: press \keys{Cmd}+\keys{Space}, type \texttt{Terminal}, press \keys{Return}.
  \item Type the line below and press \keys{Return}. A window asks to install the
        \enquote{command line developer tools}: click \menu{Install}, then \menu{Agree}. It
        downloads for 5--20 minutes. (If it says the tools are already installed, great.)
\end{checklist}
\begin{shellcode}
xcode-select --install
\end{shellcode}
\begin{checklist}
  \item Check it worked: \texttt{git --version} should print something like
        \texttt{git version 2.39.5 (Apple Git-154)}.
\end{checklist}

\subsection{5.2 git identity: a nickname, never a real name}
Git writes a name and an email into every save point. We use a \textbf{nickname} and a clearly
fake \textbf{placeholder email} (\texttt{example.com} is reserved for examples and never receives
mail). Type, one line at a time, on each Mac (replace \texttt{Pixel Fox} with the student's or
pair's chosen nickname):
\begin{shellcode}
git config --global user.name "Pixel Fox"
git config --global user.email "pixelfox@example.com"
git config --global init.defaultBranch main
git config --global core.editor nano
\end{shellcode}
\begin{checklist}
  \item \texttt{git config --global --list} shows the nickname and the placeholder email.
  \item The editor line matters: if a student forgets \texttt{-m "message"}, git opens the
        friendly \texttt{nano} editor instead of the confusing \texttt{vim} (see Troubleshooting).
  \item Two students sharing one Mac login? Run the \texttt{user.name} line \emph{without}
        \texttt{--global} inside each student's own project folder after \texttt{git init}.
\end{checklist}

\subsection{5.3 Thonny and pytest}
\begin{checklist}
  \item Thonny 4 or newer (\menu{Thonny} $\rightarrow$ \menu{About Thonny}) with Python 3.10 or newer.
  \item pytest is installed (it was in Level 2): \menu{Tools} $\rightarrow$ \menu{Manage packages\ldots},
        search \texttt{pytest}, click \menu{Install} if it is not listed as installed.
  \item Find \menu{Tools} $\rightarrow$ \menu{Open system shell\ldots}. This opens a Terminal that
        uses \emph{Thonny's} Python, which has pytest. \textbf{Run tests from this Terminal.}
        (Plain Terminal's \texttt{python3} is Apple's own copy, which does not have pytest.)
  \item Test it: in that shell type \texttt{python3 -m pytest --version}. It prints a version number.
\end{checklist}

\subsection{5.4 Copy the course files}
\begin{checklist}
  \item On each Mac make a folder \texttt{Documents/coding-L3-<nickname>} (one per student, or per pair).
  \item Copy \texttt{3-advanced/exercises} and \texttt{3-advanced/project/starter} into it.
  \item Students create their own git repositories \emph{inside their copy} (week 1 shows how).
        Never let students run git commands inside the curriculum folder you received: it may
        already be a git repository, and a repository inside a repository is confusing.
  \item Keep \texttt{project/reference} (the Maze Muncher solution) on your Mac only until week 12.
\end{checklist}

\subsection{5.5 Printing}
\begin{tabularx}{\linewidth}{@{}>{\raggedright\arraybackslash}p{4.8cm}X@{}}
\toprule
\rowcolor{kcaccent!15} When & Print (one per student unless noted)\\ \midrule
Before week 1 & Placement quiz (from \texttt{assessments.pdf}); cheat sheet (double-sided); glossary;
progress tracker; parent notes (one per family)\\
Every week & Workbook chapter and homework set for that week\\
Week 6 & Mid-level checkpoint\\
Week 9 & Project brief\\
Week 12 & Exit assessment; certificates (write names by hand)\\
For you only & Lesson plans, answer key, this guide, project teacher guide (keep in a folder
students cannot see)\\
\bottomrule
\end{tabularx}

\subsection{5.6 Accessibility}
\begin{checklist}
  \item \textbf{Dyslexia-friendly edition:} every PDF also exists in a Lexend, larger-spaced,
        cream-page edition (\texttt{build/pdf/dyslexia/}). Offer it to everyone; some students
        simply prefer it.
  \item \textbf{Thonny font size:} \menu{View} $\rightarrow$ \menu{Increase font size}
        (\keys{Cmd}+\keys{+}), or \menu{Tools} $\rightarrow$ \menu{Options\ldots} $\rightarrow$
        \menu{Theme \& Font}.
  \item \textbf{Terminal font size:} \keys{Cmd}+\keys{+} in the Terminal window; to keep it,
        \menu{Terminal} $\rightarrow$ \menu{Settings} $\rightarrow$ \menu{Profiles} $\rightarrow$ \menu{Text}.
  \item \textbf{Color filters:} \menu{System Settings} $\rightarrow$ \menu{Accessibility}
        $\rightarrow$ \menu{Display} $\rightarrow$ \menu{Color filters} (helpful for color
        blindness). All course colors are from the color-blind-safe Okabe--Ito palette, and games
        always use shapes and letters too.
  \item \textbf{Zoom:} \menu{Accessibility} $\rightarrow$ \menu{Zoom} lets students enlarge any
        part of the screen during demos.
\end{checklist}

%% =====================================================================
\section{6. The 60-minute session}

Every lesson plan follows the same shape, so after two weeks it feels routine.

\begin{timeplan}
  \timeslot{0--8}{Warm-up}{An unplugged activity or a quick puzzle that sets up today's idea. Students
  open Thonny and run \texttt{git log --oneline} in their project folder to see where they left off.}
  \timeslot{8--18}{Concept}{Read the teacher script and the analogy. Draw the diagram from the
  workbook on paper. Ask the check question before moving on.}
  \timeslot{18--28}{Guided demo}{You (or a confident student) type the worked example while everyone
  watches and predicts. Make one deliberate mistake and let them spot it.}
  \timeslot{28--52}{Hands-on build}{Pairs work through the workbook exercises: driver/navigator, swap
  at the timer every 8 minutes; check with \texttt{check\_my\_work.py}; commit when something works;
  the last 5--6 minutes are creative free-build.}
  \timeslot{52--60}{Wrap-up}{Two-minute show-and-tell, one reflection prompt, exit ticket, final
  commit (\texttt{git commit -m "..."}), hand out homework.}
\end{timeplan}

\begin{tip}[Your 10 minutes of prep]
(1) Read the lesson plan's at-a-glance table and teacher scripts. (2) Open the week's exercise
folder and run one solution. (3) Print the workbook chapter and homework. (4) Put the materials
(cards, sticky notes, timer) on the table. That is all.
\end{tip}

%% =====================================================================
\section{7. Collaboration at Level 3}

\subsection{Three roles: driver, navigator, reviewer}
\begin{pairroles}[Roles at Level 3]
\textbf{Driver} types and says out loud what they are doing. \textbf{Navigator} reads the task,
thinks one step ahead and spots bugs. \textbf{Reviewer} (new at Level 3) reads finished code with
the checklist below before it is committed. Swap driver and navigator every 8 minutes; the pair
becomes reviewer for the other pair in the last 5 minutes of the build.

\textbf{With 1 student:} you are the navigator and the reviewer. \textbf{With 3:} rotate driver,
navigator, reviewer. \textbf{With 4:} two pairs, who review each other's code.
\end{pairroles}

\subsection{Code review checklist (printed on the cheat sheet)}
\begin{checklist}
  \item It runs without errors, and the tests pass.
  \item Names say what things are (\texttt{player\_score}, not \texttt{x}).
  \item Each function or method does one job.
  \item Comments explain \emph{why}, not just what.
  \item No real names, addresses, passwords or photos anywhere.
  \item The commit message says what changed, starting with a verb.
  \item Feedback is kind and useful: \enquote{two stars and a wish}.
\end{checklist}

\subsection{git as a team tool}
Keep git simple: \textbf{one repository per pair, one branch, no merging}.
\begin{itemize}[itemsep=2pt]
  \item The repository lives on the Mac the pair uses. Whoever is driving makes the commit; the
        message can name both (\texttt{"Add Knight class (Ada + Kai)"}).
  \item Commit every time something works and at the end of every session. Small commits with
        clear messages are the habit we are building; the project needs 8 or more.
  \item To move a project to another Mac, copy the whole folder (USB stick or AirDrop). The hidden
        \texttt{.git} folder inside travels with it, so the history comes too.
  \item Branches, merging and online services such as GitHub are deliberately left out. They are
        great later, but they add confusion without adding much for a home group.
\end{itemize}

%% =====================================================================
\section{8. The weekly threads}

\begin{description}[leftmargin=0pt,itemsep=4pt,font=\bfseries]
\item[Bug hunts.] Every week has a short broken program to fix. Use the routine \enquote{Read the
error from the bottom up $\rightarrow$ find the line $\rightarrow$ guess $\rightarrow$ test one
change}. From week 6 students also use Thonny's debugger (the bug button, then \menu{Step over}).
Celebrate bugs found: \enquote{Nice catch!}
\item[Unplugged.] Paper and body activities make abstract ideas visible: save-point timeline
cards, cookie-cutter classes, human sorting race, plate stack and lunch queue, Russian dolls,
paper survey, \enquote{train the teacher}, CRC cards. Do not skip them; they are where most
\enquote{aha} moments happen.
\item[Free-build.] Five minutes at the end of the build for students to make something their own
(new colors, a new character, a silly feature). It keeps motivation high, especially for the
artists in the group.
\item[Reflection.] One prompt at the end of each session, answered aloud or in the workbook.
Rotate who answers first. Good follow-up: \enquote{What would you tell someone starting this
tomorrow?}
\end{description}

%% =====================================================================
\section{9. Assessment plan}

Assessment at home is for \emph{learning}, not for grades: it tells you what to revisit and who
needs a different path. Papers are in \texttt{assessments.pdf}; answers, rubrics and scoring are
in the answer key (\enquote{Assessments} section).

\begin{longtable}{@{}>{\bfseries}p{2.9cm}>{\raggedright\arraybackslash}p{2.2cm}p{\dimexpr\linewidth-5.1cm-4\tabcolsep}@{}}
\toprule
\rowcolor{kcaccent!15} What & When & How you use it\\ \midrule\endhead
\bottomrule\endfoot
Placement quiz & Before week 1 (ideally at home the week before), 15 min & Covers Level 2:
functions with return, lists, dicts, \texttt{while}, files, a pytest test. Score out of 20.
\textbf{16--20:} standard path, offer extension tasks. \textbf{10--15:} standard path with the
support notes in weeks 1--3, plus the refresher listed in the answer key. \textbf{Below 10:} run a
refresher session first (Level 2 weeks 2, 3, 5 and 7 homework), and pair this student with you
or a confident partner. Retake the missed questions after the refresher.\\ \midrule
Formative checks & Every session & The lesson's exit ticket (one question, answer in the plan),
the self-check boxes in homework and watching the hands-on build. Note anything worrying on the
progress tracker.\\ \midrule
Mid-level checkpoint & Week 6, 15 min & Classes, inheritance, search/sort tracing, stacks and
queues. Use the results to choose the focus of the week 6 review and whether to use buffer B1
before week 7.\\ \midrule
Exit assessment & Week 12, 20 min & Recursion, grids, data, AI concepts: read, predict, fix,
design and explain. Combined with the project, it tells you whether a student is ready to work
independently on Python projects.\\ \midrule
Project rubric & Weeks 9--12 & Six criteria: creativity, working code, design/OOP, testing, git
history, explaining design decisions. Explaining counts as much as code. Full rubric in the
project teacher guide.\\ \midrule
Progress tracker & Every week & One page per student (\texttt{progress-tracker.pdf}): check the
week's \enquote{I can} statements together during wrap-up.\\
\end{longtable}

\begin{teachernote}[Every student gets a voice]
Each paper ends with student-friendly \enquote{I can\ldots} statements scored
\textbf{Just starting / Getting there / Got it / Expert}. Ask students to fill these in
\emph{before} you score, then compare. A student who rates themself \enquote{Got it} but struggled
is a conversation, not a correction.
\end{teachernote}

%% =====================================================================
\section{10. Differentiation}

Every lesson plan has a support path and an extension path. In short:

\begin{longtable}{@{}>{\bfseries}p{0.8cm}>{\raggedright\arraybackslash}p{7cm}>{\raggedright\arraybackslash}p{\dimexpr\linewidth-7.8cm-4\tabcolsep}@{}}
\toprule
\rowcolor{kcaccent!15} Wk & Support path (needs more time) & Extension path (finished early)\\ \midrule\endhead
\bottomrule\endfoot
1 & Only \texttt{init}, \texttt{add}, \texttt{commit}, \texttt{log}; parent types first, student copies & Use \texttt{git diff} and \texttt{git restore} to undo a deliberate mistake\\ \midrule
2 & Fill in a half-written class; one attribute, one method & Add a \texttt{\_\_str\_\_} method and a list of pets\\ \midrule
3 & One parent class, one child that overrides one method & Composition: a Team that has a list of characters\\ \midrule
4 & Linear search only, count steps on paper & Binary search on 1000 numbers, compare step counts in a table\\ \midrule
5 & Use \texttt{sorted(key=...)} instead of writing a sort & Write bubble sort and count the swaps; undo with a stack\\ \midrule
6 & Redo one favorite exercise from weeks 2--5 & Fix the hardest bug-hunt station with the debugger only\\ \midrule
7 & Draw fractal trees by changing depth/angle only & Flood fill on a grid; then try BFS shortest path\\ \midrule
8 & Count with a dict; chart already drawn & Clean a messier CSV; add averages and a second chart\\ \midrule
9 & Classify by hand first, then run the given classifier & k-nearest neighbors (3 vote); compare accuracy\\ \midrule
10 & Start from a starter template; two classes only & Design an extra class with composition; write tests first\\ \midrule
11 & Five simple tests (one per method); README from the template & Edge-case tests; tests for the saved file with \texttt{tmp\_path}\\ \midrule
12 & Present with a prompt card; demo one feature & Explain one alternative design and why they rejected it\\
\end{longtable}

%% =====================================================================
\section{11. Digital citizenship and AI ethics}

One short conversation every week (the \enquote{Digital citizenship \& safety} box in the plan).
Keep it a discussion, not a lecture: ask, listen, then share one fact.

\begin{longtable}{@{}>{\bfseries}p{0.8cm}>{\raggedright\arraybackslash}p{4.6cm}p{\dimexpr\linewidth-5.4cm-4\tabcolsep}@{}}
\toprule
\rowcolor{kcaccent!15} Wk & Topic & A question to ask\\ \midrule\endhead
\bottomrule\endfoot
1 & Never commit secrets & \enquote{Git remembers everything forever. What should never go into a commit?}\\ \midrule
2 & Kind virtual pets & \enquote{Some games make pets sad if you do not play every day. Is that fair?}\\ \midrule
3 & Representation & \enquote{Who gets to be the hero in games you play? Who is missing?}\\ \midrule
4 & Search engines and filter bubbles & \enquote{Why might two people get different search results?}\\ \midrule
5 & Fair, kind leaderboards & \enquote{How could a leaderboard make someone feel bad? How could we fix that?}\\ \midrule
6 & Growth mindset & \enquote{What is something that was hard in week 2 and is easier now?}\\ \midrule
7 & Credit for algorithms and art & \enquote{If you use someone's fractal idea, how do you give credit?}\\ \midrule
8 & Consent and anonymous data & \enquote{May we put someone's survey answer in our chart without asking?}\\ \midrule
9 & AI ethics: chatbots, deepfakes, bias & \enquote{What should you never type into a chatbot? How could an AI be unfair?}\\ \midrule
10 & Licenses & \enquote{Can you use any picture or sound you find? What does a license say?}\\ \midrule
11 & Accessibility and usability & \enquote{Could someone who is color-blind play your game?}\\ \midrule
12 & Showcase safety; next steps online & \enquote{If you join a coding website later, what will you keep private?}\\
\end{longtable}

\begin{safety}[The AI-chatbot rule for this course]
Students do not need any AI chatbot for this course. If your family uses one at home: never type
real names, addresses, school names, photos or passwords into it; check its answers (it can be
confidently wrong); and say when you used it, the same way you credit any helper. Many chatbot
services also have minimum ages (often 13); please follow them.
\end{safety}

%% =====================================================================
\section{12. Troubleshooting}

\subsection{git}
\begin{longtable}{@{}>{\raggedright\arraybackslash}p{5.2cm}p{\dimexpr\linewidth-5.2cm-2\tabcolsep}@{}}
\toprule
\rowcolor{kcaccent!15} What you see & What to do\\ \midrule\endhead
\bottomrule\endfoot
\texttt{fatal: not a git repository (or any of the parent directories): .git} & The Terminal is
in the wrong folder. Type \texttt{pwd} to see where you are, then \texttt{cd} into the project
folder (tip: type \texttt{cd }, then drag the folder from Finder into the Terminal window and
press \keys{Return}). If the folder was never set up, run \texttt{git init} there.\\ \midrule
\texttt{Author identity unknown \ldots Please tell me who you are} & Do step 5.2 (nickname and
placeholder email), then commit again.\\ \midrule
\texttt{nothing to commit, working tree clean} & Not an error. Either everything is already
committed, or the file was not saved in Thonny yet: press \keys{Cmd}+\keys{S}, then commit.\\ \midrule
\texttt{nothing added to commit but untracked files present} & Forgot \texttt{git add}. Run
\texttt{git add .} then commit again.\\ \midrule
A screen full of \texttt{\textasciitilde} signs, typing does nothing sensible & Git opened the
\texttt{vim} editor because the \texttt{-m "message"} part was missing. Press \keys{Esc}, type
\texttt{:q!} and press \keys{Return}. Then commit with \texttt{-m}. (In \texttt{nano}: press
\keys{Ctrl}+\keys{X}.) Step 5.2 sets \texttt{nano} to avoid this.\\ \midrule
\texttt{You are in 'detached HEAD' state} & Someone used \texttt{git checkout <code>} to look at
an old save point. Nothing is lost. Type \texttt{git switch main} to go back to the latest work.
To get one old file back instead, use \texttt{git restore --source <code> file.py}.\\ \midrule
\texttt{CONFLICT} or \texttt{merge} messages & Only happens if someone made a branch. Type
\texttt{git merge --abort}, then \texttt{git switch main}. Our single-branch workflow avoids
merges completely.\\ \midrule
\enquote{\texttt{git restore} deleted my changes!} & \texttt{restore} throws away changes that
were never committed; it cannot be undone. Turn it into a lesson: this is why we commit often.\\ \midrule
\texttt{xcrun: error: invalid active developer path} & Usually after a macOS update. Run
\texttt{xcode-select --install} again.\\
\end{longtable}

\subsection{Python, data and turtle}
\begin{longtable}{@{}>{\raggedright\arraybackslash}p{5.2cm}p{\dimexpr\linewidth-5.2cm-2\tabcolsep}@{}}
\toprule
\rowcolor{kcaccent!15} What you see & What to do\\ \midrule\endhead
\bottomrule\endfoot
\texttt{RecursionError: maximum recursion depth exceeded} & The function keeps calling itself.
Check: (1) is there a base case (\texttt{if n == 0: return})? (2) does every call make the problem
\emph{smaller} (\texttt{n - 1}, \texttt{length * 0.7})? (3) flood fill: mark a square as visited
\emph{before} recursing, and keep grids small (about 30 $\times$ 30 or less).\\ \midrule
\texttt{UnicodeDecodeError} when opening a CSV, or a \texttt{KeyError} for the first column
(e.g. \texttt{'\textbackslash ufeffname'}) & The file came from a spreadsheet app with a hidden
marker at the start. Open it with \texttt{encoding="utf-8-sig"}, or re-export it from Numbers:
\menu{File} $\rightarrow$ \menu{Export To} $\rightarrow$ \menu{CSV\ldots}, text encoding
\menu{Unicode (UTF-8)}.\\ \midrule
\texttt{TypeError: can only concatenate str (not "int")} with CSV data & Everything read from a
CSV is text. Convert with \texttt{int(row["points"])}; skip empty rows first.\\ \midrule
Blank lines between rows when writing a CSV & Open the file with \texttt{newline=""}.\\ \midrule
Turtle window does not respond to keys & Click on the turtle window first; make sure the program
calls \texttt{screen.listen()}; check the key name (\texttt{"Up"}, \texttt{"space"}, \texttt{"a"}).\\ \midrule
Turtle drawing is very slow & Add \texttt{screen.tracer(0)} at the start and
\texttt{screen.update()} after drawing each frame.\\ \midrule
\texttt{turtle.Terminator} after closing a window & Harmless. Press Run again (or \menu{Run}
$\rightarrow$ \menu{Stop/Restart backend} first).\\ \midrule
Window freezes / spinning beach ball & An endless \texttt{while True} loop without \texttt{ontimer},
or a missing base case. Press the red \menu{Stop} button in Thonny.\\
\end{longtable}

\subsection{pytest}
\begin{longtable}{@{}>{\raggedright\arraybackslash}p{5.2cm}p{\dimexpr\linewidth-5.2cm-2\tabcolsep}@{}}
\toprule
\rowcolor{kcaccent!15} What you see & What to do\\ \midrule\endhead
\bottomrule\endfoot
\texttt{No module named pytest} & You are in plain Terminal. Use Thonny's \menu{Tools}
$\rightarrow$ \menu{Open system shell\ldots} (step 5.3), or install pytest with
\menu{Tools} $\rightarrow$ \menu{Manage packages\ldots}.\\ \midrule
\texttt{collected 0 items} / \texttt{no tests ran} & Test files must be named
\texttt{test\_something.py} and test functions must start with \texttt{test\_}. Also check you
\texttt{cd}-ed into the project folder.\\ \midrule
\texttt{ModuleNotFoundError: No module named 'dragon\_maze'} & The program was renamed or the
test's \texttt{import} line does not match the file name. File names must not contain spaces or
dashes to be imported.\\ \midrule
Tests open a turtle window or hang & The test imported code that draws at import time. Keep the
drawing inside functions and the main program under \texttt{if \_\_name\_\_ == "\_\_main\_\_":}.
Test the logic, not the drawing.\\
\end{longtable}

%% =====================================================================
\section{13. Pacing, buffers and holidays}

\subsection{Pacing advice}
\begin{itemize}[itemsep=2pt]
  \item 12 core sessions + 2 buffer sessions = 14 sessions. Plan them across \textbf{14--16
        calendar weeks} so that holidays and sick days do not cause stress.
  \item \textbf{Never skip weeks 6 and 12.} The checkpoint tells you what to fix before the second
        half; the showcase is the reward the whole level builds towards.
  \item If a lesson runs over, stop at the wrap-up anyway: the reflection and final commit are
        more valuable than finishing the last exercise. Unfinished exercises become homework.
  \item Use \textbf{B1 (Algorithm Olympics)} after week 6 if the checkpoint shows gaps in
        searching/sorting, or any week energy is low. Use \textbf{B2 (Mini Game Jam)} before the
        project (week 9) or around a holiday when attendance is patchy.
\end{itemize}

\subsection{If you lose sessions}
\begin{tabularx}{\linewidth}{@{}>{\bfseries}p{2.5cm}X@{}}
\toprule
\rowcolor{kcaccent!15} Sessions lost & What to do\\ \midrule
1 & Drop B2. If B2 is already used, drop B1.\\
2 & Drop both buffers.\\
3 & Drop both buffers and merge weeks 4 and 5: binary search only (skip counting steps on paper),
then \texttt{sorted(key=...)} and stacks/queues; set bubble sort as homework.\\
4 or more & Also merge weeks 7 and 8 (fractal tree + one CSV chart, skip flood fill), or move the
project milestones M3/M4 into one session and let students finish the README at home. Keep
weeks 6, 9, 10 and 12 intact.\\
\bottomrule
\end{tabularx}

\subsection{Sample calendar: Saturdays, September--December 2027}
A US-friendly plan with 14 sessions in 15 calendar weeks. It skips Labor Day weekend and
Thanksgiving weekend and finishes before the winter holidays. Move it to any weekday you like.

\begin{longtable}{@{}>{\bfseries}p{2.3cm}>{\raggedright\arraybackslash}p{5.4cm}>{\raggedright\arraybackslash}p{\dimexpr\linewidth-7.7cm-4\tabcolsep}@{}}
\toprule
\rowcolor{kcaccent!15} Date & Session & Notes\\ \midrule\endhead
\bottomrule\endfoot
Sat 4 Sep & --- (no session) & Labor Day weekend. Hand out the placement quiz to do at home.\\
Sat 11 Sep & Week 1: Save Points & Collect the placement quiz; parent notes go home.\\
Sat 18 Sep & Week 2: Objects &\\
Sat 25 Sep & Week 3: Families of Objects &\\
Sat 2 Oct & Week 4: Searching &\\
Sat 9 Oct & Week 5: Sorting, Stacks, Queues & Indigenous Peoples'/Columbus Day is Mon 11 Oct.\\
Sat 16 Oct & Week 6: Review + Checkpoint & Never skip.\\
Sat 23 Oct & B1: Algorithm Olympics & Revisit gaps from the checkpoint.\\
Sat 30 Oct & Week 7: Fractal Art & Halloween weekend: spooky fractal trees welcome.\\
Sat 6 Nov & Week 8: Data Detectives & Collect the paper survey the week before.\\
Sat 13 Nov & Week 9: Machines Learn + kickoff & Hand out the project brief.\\
Sat 20 Nov & B2: Mini Game Jam & Spare session: if a week was lost, skip B2 and shift up.\\
Sat 27 Nov & --- (no session) & Thanksgiving weekend.\\
Sat 4 Dec & Week 10: Project Studio 1 &\\
Sat 11 Dec & Week 11: Project Studio 2 &\\
Sat 18 Dec & Week 12: Showcase + Graduation & Never skip. Invite family to the showcase.\\
\end{longtable}

%% =====================================================================
\section{14. Coverage}

\subsection{(a) Goals $\times$ weeks}
\textbf{I} = introduced, R = revisited, \textbf{E}\textsuperscript{+} = extended (used in a new,
harder way). Written from the master syllabus.

\begingroup\small
\setlength{\tabcolsep}{3.2pt}
\begin{tabularx}{\linewidth}{@{}>{\bfseries}l*{12}{>{\centering\arraybackslash}X}>{\centering\arraybackslash}p{0.7cm}>{\centering\arraybackslash}p{0.7cm}@{}}
\toprule
\rowcolor{kcaccent!15} Goal & 1 & 2 & 3 & 4 & 5 & 6 & 7 & 8 & 9 & 10 & 11 & 12 & B1 & B2\\ \midrule
G1 OOP     &     & \cvI & \cvE &      &      & \cvR &      &      &      & \cvE & \cvR & \cvR &      &     \\
G2 Algor.  &     &      &      & \cvI & \cvE & \cvR & \cvE &      &      &      &      & \cvR & \cvR &     \\
G3 Data str.& \cvR & \cvR &     &      & \cvI & \cvR & \cvE & \cvR &      & \cvR &      & \cvR & \cvR &     \\
G4 git     & \cvI & \cvR & \cvR & \cvR & \cvR & \cvR & \cvR & \cvR & \cvR & \cvR & \cvE & \cvR &      & \cvR\\
G5 Project &     &      &      &      &      &      &      &      & \cvI & \cvR & \cvE & \cvR &      & \cvR\\
G6 Data    &     &      &      &      &      &      &      & \cvI & \cvR &      &      & \cvR &      &     \\
G7 AI      &     &      &      &      &      &      &      &      & \cvI &      &      & \cvR &      &     \\
G8 Test/debug & \cvR & \cvR & \cvR & \cvR & \cvR & \cvI & \cvR & \cvR & \cvR & \cvR & \cvE & \cvR &  & \cvR\\
G9 Ethics  & \cvI & \cvR & \cvR & \cvR & \cvR & \cvR & \cvR & \cvR & \cvE & \cvR & \cvR & \cvR &      &     \\
\bottomrule
\end{tabularx}
\endgroup

\textbf{Notes.} G4: from week 2 every session ends with a commit (R); week 11 needs a project
history of 8+ meaningful commits (\textbf{E}\textsuperscript{+}). G8: bug hunts revisit debugging
every week; the Thonny debugger is introduced in week 6; project tests extend testing in week 11.
G3: lists and dicts are revisited in weeks 1, 2 and 8; stacks/queues are new in week 5 and grids in
week 7.

\subsection{(b) Objectives $\rightarrow$ lessons}
The table below is generated automatically from the lesson plans' objectives.

\kcinput{generated/coverage}

%% =====================================================================
\section{15. Spiral map: how big ideas come back across Levels 1--3}

\begingroup\small
\begin{longtable}{@{}>{\bfseries\raggedright\arraybackslash}p{2.6cm}>{\raggedright\arraybackslash}p{3.6cm}>{\raggedright\arraybackslash}p{4.2cm}>{\raggedright\arraybackslash}p{\dimexpr\linewidth-10.4cm-6\tabcolsep}@{}}
\toprule
\rowcolor{kcaccent!15} Big idea & Level 1 & Level 2 & Level 3\\ \midrule\endhead
\bottomrule\endfoot
Sequence \& algorithms & W1 I (unplugged), W7 R (text) & W2 R (functions), W10 E (pseudocode) & W4--5 E (named algorithms, efficiency)\\ \midrule
Loops & W2 I (Scratch), W8 R (for) & W4 E (while, nested), W9 E (game loop) & W4 R (search loops), W7 E (recursion as repetition)\\ \midrule
Events \& conditionals & W3 I, W9 R (if/elif) & W1 E (and/or), W9 E (onkey/ontimer) & W2--3 R (methods react), project E\\ \midrule
Variables \& data & W4 I, W9 R (input) & W1 E (types), W3 E (lists), W5 E (dicts), W7 E (files) & W2 E (attributes), W5 E (stacks/queues), W7 E (grids), W8 E (CSV data)\\ \midrule
Abstraction / functions & W5 I (My Blocks), W10 R (def) & W2 E (return), W8 R (testable functions) & W2--3 E (classes), W10 E (modules)\\ \midrule
Debugging \& testing & every week (bug hunt), W6 relay & W6 E (strategies), W8 E (pytest) & W6 E (debugger), project E (5+ tests)\\ \midrule
Collaboration & pair roles every week & W10 E (task split) & W1 E (git), W10 E (code review)\\ \midrule
Art \& games & W2 pen art, W4 catch game, W10 random art & W3 palettes, W4 pixel art, W9 games & W3 shape families, W7 fractals, project\\ \midrule
Citizenship \& AI ethics & privacy, credit, kindness, accessibility & data privacy, encryption, fair games, attribution & git secrets, consent, bias, AI ethics, licenses\\
\end{longtable}
\endgroup

%% =====================================================================
\section{16. Where next after Level 3}

Your graduates can now learn from real programming books on their own. Some safe next steps:

\begin{description}[leftmargin=0pt,itemsep=4pt,font=\bfseries]
\item[Books (offline, great with Thonny).] \enquote{Python for Kids} (Jason Briggs, 2nd edition);
\enquote{Mission Python} (Sean McManus); \enquote{Invent Your Own Computer Games with Python} and
\enquote{The Big Book of Small Python Projects} (Al Sweigart); later \enquote{Python Crash Course}
(Eric Matthes). Your local library probably has several.
\item[Offline projects.] Turn Maze Muncher into a level editor; build a turtle drawing program that
saves to JSON; analyze family-friendly data (weather you measure, a reading log, a sports season);
write a text adventure with classes for rooms and items; keep everything in git.
\item[Local, in-person clubs.] Library coding clubs, CoderDojo or Code Club sessions, school
robotics teams (for example FIRST LEGO League), maker spaces and summer coding camps.
\item[Online, with a parent.] Many sites (GitHub, Python forums, chat communities) require users to
be at least 13. Before then, explore together: you create and own the account, use a nickname,
share nothing personal, and agree on what can be posted. Puzzle sites with no social features are
the gentlest start.
\item[Keep the habits.] Small steps, frequent commits, tests for the tricky parts, credit for
others' work, and kindness in feedback.
\end{description}

\end{document}
